\section{Théorème directeur du réseau radical et de la loi puissance--lacune}

L’apparition de \(369\) et de \(693\) ne se réduit pas à une ressemblance graphique. Deux mécanismes exacts convergent :

\begin{enumerate}
\item \textbf{mécanisme algébrique nonadique :} \(3,6,9\) sont les représentants visuels du radical nilpotent \(\mathfrak m=(3)\subset\mathbb Z/9\mathbb Z\) ; les mots \(369\), \(693\) et \(936\) sont les trois phases de son orbite multiplicative ;
\item \textbf{mécanisme archimédien des lacunes :} la différence entre la longueur décimale de \(9^n\) et la frontière \(10^n\) est gouvernée exactement par la même rotation irrationnelle \(\delta_{\rm vac}=\log_{10}(10/9)\) qui organise déjà le calendrier des lacunes.
\end{enumerate}

Pour \(n=9^9\), les deux mécanismes se rencontrent dans l’identité

\[
9^9-\left\lceil 9^9\log_{10}\frac{10}{9}\right\rceil
=369\,693\,099,
\]

de sorte que

\[
\#\operatorname{cifras}\bigl(9^{9^9}\bigr)=369\,693\,100.
\]

Le bloc \(369\mid693\) est donc une sortie exacte du lecteur logarithmique de la puissance nonadique ; ce n’est pas le préfixe décimal de \(9^{9^9}\). L’unité finale provient de l’opérateur universel qui transforme l’ordre décimal en longueur décimale.

Cette identité ouvre une lecture structurelle de plus grande portée : somme, produit, puissance, soustraction, division, racine et logarithme peuvent être organisés comme une hiérarchie de transports et de quotients. Chaque montée d’opération compactifie des histoires ; le TPK conserve les coordonnées qui permettent de les reconstruire.


\section{Structure radicale de l’APP sur \(\mathbb Z/9\mathbb Z\)}

\subsection{Anneau local nonadique \(\mathbb Z/9\mathbb Z\) et radical nilpotent}

Soit

\[
R=\mathbb Z/9\mathbb Z.
\]

Son groupe des unités et son idéal maximal sont

\[
R^\times=\{1,2,4,5,7,8\},
\qquad
\mathfrak m=(3)=\{0,3,6\}=\{9,3,6\}_{\rm vis}.
\]

De plus,

\[
\mathfrak m^2=0.
\]

L’égalité précédente est une égalité d’idéaux : tout produit de deux éléments de \(\mathfrak m\) s’annule modulo neuf. Elle ne doit pas être confondue avec le produit cartésien \(\mathfrak m\times\mathfrak m\), qui contient neuf positions de l’APP.

\subsection{Suite exacte du multiplicateur \(3\)}

La multiplication par trois,

\[
M_3:R\longrightarrow R,
\qquad x\longmapsto3x,
\]

a

\[
\ker M_3=\mathfrak m,
\qquad
\operatorname{im}M_3=\mathfrak m.
\]

Elle induit donc un isomorphisme d’espaces sur \(\mathbb F_3\),

\[
\overline M_3:R/\mathfrak m\xrightarrow{\sim}\mathfrak m,
\qquad
\overline x\longmapsto3x.
\]

En représentants visuels,

\[
\overline0\longmapsto9,
\qquad
\overline1\longmapsto3,
\qquad
\overline2\longmapsto6.
\]

Ainsi, la bande \(3,6,9\) n’est pas un ornement décimal. C’est une copie interne du quotient ternaire logée dans le radical nonadique.

\subsection{Orbite cyclique \(369\to693\to936\to369\)}

La ligne de la table multiplicative correspondant à \(3\) est

\[
3,6,9,3,6,9,3,6,9.
\]

Le changement de phase \(j\mapsto j+1\) produit

\[
369\longmapsto693\longmapsto936\longmapsto369.
\]

Ces trois mots sont les trois lectures cycliques de l’application

\[
\mathbb F_3\longrightarrow\mathfrak m,
\qquad
\bar j\longmapsto3j.
\]

La ligne \(6\), en satisfaisant \(6\equiv-3\pmod9\), inverse l’orientation du même cycle. La marque \(9\) remplit deux fonctions typées :

\[
9\bmod9=0,
\qquad
\operatorname{dr}_9(9)=9.
\]

Le module publie la classe nulle ; la racine numérique positive conserve la marque de clôture. C’est précisément cette différence de lecteurs qui permet à \(9\) d’être à la fois zéro résiduel et clôture visible.

\subsection{Croix radicale de l’APP et réduction coordonnée à \(\mathbb F_3^2\)}

Sur

\[
X_9=R^2,
\]

l’union

\[
\mathscr C_{369}
=(\mathfrak m\times R)\cup(R\times\mathfrak m)
\]

est la croix formée par les lignes et les colonnes \(3,6,9\). Elle contient

\[
3\cdot9+3\cdot9-3\cdot3=45
\]

positions. Son intersection centrale \(\mathfrak m\times\mathfrak m\) contient neuf positions. Dans l’observable multiplicative, la somme de la croix est

\[
297=216+81=3(72+27).
\]

La réduction coordonnée produit

\[
X_9/(\mathfrak m\times\mathfrak m)\cong\mathbb F_3^2
\]

comme ensemble quotient typé par les classes modulo trois. Le relèvement inverse exige de conserver la coordonnée à l’intérieur de chaque fibre de trois éléments.

\subsection{Noyau \(7227\), décomposition spectrale et lien avec la croix radicale}

Le bloc central de l’APP est

\[
B_c=
\begin{pmatrix}
7&2\\
2&7
\end{pmatrix}
=7I+2R
=9P_++5P_-.
\]

Sa lecture par lignes produit \(72\mid27\) ; sa lecture par diagonales produit \(77\mid22\). On vérifie

\[
72+27=77+22=99,
\]

\[
72-27=45=\det B_c,
\]

\[
77-22=55=54+1.
\]

La relation avec la croix radicale n’est pas nominale :

\[
3\cdot72=216,
\qquad
3\cdot27=81,
\qquad
3\cdot99=297.
\]

Le bloc unitaire central et le réseau non unitaire \(3,6,9\) sont ainsi reliés par une identité locale–globale exacte. La hiérarchie de la différence somme–produit conserve ses types :

\[
4\quad\text{(saut spectral local)},\qquad
2\quad\text{(couplage hors diagonale)},
\]

\[
6\quad\text{(différence comprimée modulo trois)},\qquad
54\quad\text{(déploiement bidimensionnel)}.
\]


\section{Loi puissance--lacune pour les échelles \texorpdfstring{\(9^n\)}{9 à la puissance n} et \texorpdfstring{\(10^n\)}{10 à la puissance n}}

\subsection{Fonctions de longueur décimale et de lacune cumulée}

Pour \(n\ge1\), soit

\[
D_9(n)=\#\operatorname{cifras}(9^n)
=\left\lfloor n\log_{10}9\right\rfloor+1.
\]

La frontière décimale \(10^n\) possède \(n+1\) chiffres. Nous définissons la lacune cumulée de la puissance nonadique par

\[
V_9(n)=(n+1)-D_9(n).
\]

Nous introduisons les deux nombres complémentaires

\[
\rho_{\rm vac}=\log_{10}9,
\qquad
\delta_{\rm vac}=1-\rho_{\rm vac}
=\log_{10}\frac{10}{9}.
\]

\subsection{Identité exacte entre longueur décimale et rotation des lacunes}

\textbf{Théorème.} Pour tout entier \(n\ge1\),

\[
\boxed{V_9(n)=\left\lceil n\delta_{\rm vac}\right\rceil},
\]

et, par conséquent,

\[
\boxed{D_9(n)=n-\left\lceil n\delta_{\rm vac}\right\rceil+1}.
\]

\textbf{Démonstration.} Comme \(\delta_{\rm vac}\) est irrationnel, \(n\delta_{\rm vac}\notin\mathbb Z\). Alors

\[
\begin{aligned}
D_9(n)
&=\lfloor n(1-\delta_{\rm vac})\rfloor+1\\
&=\lfloor n-n\delta_{\rm vac}\rfloor+1\\
&=n-\lceil n\delta_{\rm vac}\rceil+1.
\end{aligned}
\]

En soustrayant de \(n+1\), on obtient la première identité. \(\square\)

\subsection{Encodage mécanique des lacunes par les puissances nonadiques}

Les accroissements de la lacune sont

\[
Z_n=V_9(n+1)-V_9(n)
=\lceil(n+1)\delta_{\rm vac}\rceil
-\lceil n\delta_{\rm vac}\rceil
\in\{0,1\}.
\]

C’est exactement le mot mécanique de pente \(\delta_{\rm vac}\) utilisé par le calendrier HMT des lacunes. Ce calendrier admet donc une lecture supplémentaire entièrement exacte :

\begin{quote}
\(Z_n=1\) si et seulement si, en passant de \(9^n\) à \(9^{n+1}\), on perd un nouveau rang décimal par rapport à la frontière \(10^n\to10^{n+1}\).
\end{quote}

La puissance n’ajoute pas une coïncidence à la théorie des lacunes : elle fournit sa réalisation naturelle comme compteur de perte de rang décimal.

\subsection{Loi \(21/22\) comme cas du théorème puissance--lacune}

Puisque

\[
\delta_{\rm vac}^{-1}
=21.85434532678283256\ldots,
\]

les indices pour lesquels \(Z_n=1\) sont séparés par \(21\) ou \(22\). Les vingt et un premiers exposants satisfont

\[
D_9(n)=n,\qquad1\le n\le21,
\]

tandis que

\[
D_9(22)=21.
\]

En particulier,

\[
9^9=387\,420\,489
\]

possède exactement neuf chiffres. Cette conservation initiale du rang n’est pas une règle indéfinie : le calendrier des lacunes détermine exactement quand elle cesse d’être vérifiée.

\subsection{Évaluation exacte en \(n=9^9\) et longueur de \(9^{9^9}\)}

Soit

\[
n_*=9^9=387\,420\,489.
\]

Alors

\[
n_*\delta_{\rm vac}
=17\,727\,389.3684296412564569\ldots,
\]

et, par conséquent,

\[
\boxed{V_9(n_*)=17\,727\,390}.
\]

La coordonnée conservée après déduction des lacunes est

\[
n_*-V_9(n_*)
=387\,420\,489-17\,727\,390
=369\,693\,099.
\]

Par conséquent,

\[
\boxed{\operatorname{ord}_{10}\bigl(9^{9^9}\bigr)=369\,693\,099},
\]

\[
\boxed{D_9(9^9)=369\,693\,100}.
\]

On obtient en outre les congruences

\[
9^9\equiv0\pmod9,\qquad
V_9(9^9)\equiv0\pmod9,
\]

\[
369\,693\,099\equiv0\pmod9,\qquad
369\,693\,100\equiv1\pmod9.
\]

Les sommes de chiffres correspondantes sont

\[
s_{10}(9^9)=45,\qquad
s_{10}(V_9(9^9))=36,
\]

\[
s_{10}(369\,693\,099)=54,\qquad
s_{10}(369\,693\,100)=37.
\]

Le retour \(9\to1\) apparaît ici avec des types précis :

\begin{enumerate}
\item \(9^9\) et la lacune cumulée appartiennent à la classe nulle modulo neuf ;
\item leur différence, l’ordre décimal, reste dans la classe nulle et possède une somme de chiffres \(54\) ;
\item l’opérateur universel \(D=\operatorname{ord}+1\) ajoute l’unité de clôture et publie la classe positive \(1\).
\end{enumerate}

\subsection{Domaine de validité de la loi puissance--lacune}

Il est démontré :

\begin{itemize}
\item le théorème puissance–lacune pour tout \(n\) ;
\item l’identité avec le mot mécanique des lacunes ;
\item l’évaluation exacte en \(n=9^9\) ;
\item l’ordre et la longueur décimale de \(9^{9^9}\) ;
\item les congruences et sommes de chiffres précédentes.
\end{itemize}

On n’affirme pas que \(9\) soit l’unique base ayant une propriété analogue dans tout domaine, que chaque apparition externe de \(369\) provienne de cette tour, ni que la tour remplace la construction des neuf phases ou le certificat de génération de \(\pi,e,\varphi\).


\section{Hiérarchie opératorielle, réversibilité et mémoire généalogique}

\subsection{Translation, dilatation et exponentiation comme niveaux opératoriels}

Les opérations élémentaires peuvent être vues comme des actions :

\[
T_a(x)=x+a,\qquad
M_a(x)=ax,\qquad
P_k(x)=x^k.
\]

Leurs inverses partielles sont, respectivement, la soustraction, la division et l’extraction de racines. Le passage d’une opération à la suivante modifie la structure des fibres.

\subsection{Réversibilité stratifiée dans \(\mathbb Z/9\mathbb Z\)}

\textbf{Proposition.} Dans l’anneau nonadique :

\begin{enumerate}
\item toute translation \(T_a\) est bijective ;
\item \(M_a\) est bijective si et seulement si \(a\in R^\times\) ;
\item si \(a\in\mathfrak m\), la multiplication possède des fibres non triviales et son image est contenue dans \(\mathfrak m\) ;
\item pour \(x\in\mathfrak m\), \(x^k=0\) pour tout \(k\ge2\) ;
\item sur \(R^\times\cong C_6\), l’application \(P_k\) est un automorphisme si et seulement si \(\gcd(k,6)=1\).
\end{enumerate}

\textbf{Démonstration.} La première affirmation utilise la structure de groupe additif. La deuxième est la caractérisation des unités. La troisième découle de \(aR\subseteq\mathfrak m\). La quatrième procède de \(\mathfrak m^2=0\). La cinquième est la condition pour que la multiplication par \(k\) soit un automorphisme du groupe cyclique d’ordre six. \(\square\)

Le résultat produit une hiérarchie exacte :

\[
\text{somme : réversibilité globale},
\]

\[
\text{produit : réversibilité restreinte aux unités},
\]

\[
\text{puissance : dynamique cyclique sur les unités et effondrement nilpotent dans }\mathfrak m.
\]

Les racines ne sont des fonctions que lorsque la puissance correspondante est injective sur le domaine choisi ; sur le domaine total, ce sont des correspondances à plusieurs branches ou des relations vides.

\subsection{Lecture du Topological Phase Kernel sur la hiérarchie opératorielle}

Cette stratification suggère une définition fonctionnelle :

\begin{quote}
Le TPK n’ajoute pas une opération extérieure à la somme, au produit et à la puissance. Il conserve le feuillet, le quotient, la retenue, l’orientation, la route et la frontière que chaque évaluation ordinaire élimine lorsqu’elle identifie plusieurs histoires à une même sortie.
\end{quote}

La somme et le produit sont les deux premières évaluations sur l’APP. La puissance est une composition répétée du produit ; la racine sélectionne ou reconstruit des branches de cette composition. Le registre distingue la valeur publiée de l’histoire opératorielle qui la produit.

\subsection{Descente logarithmique entre niveaux consécutifs de la hiérarchie}

Dans un domaine positif,

\[
\log(xy)=\log x+\log y,\qquad
\log(x^k)=k\log x,\qquad
\log(a^x)=x\log a.
\]

Le logarithme transforme le produit en somme, la puissance en produit et l’exponentielle en dilatation. Le défaut apparaît lorsque le lecteur ultérieur discrétise :

\[
\mathcal L_{10}(a^n)
=\lfloor n\log_{10}a\rfloor+1.
\]

La partie entière introduit le calendrier des franchissements de frontière. Pour \(a=9\), ce calendrier est précisément le mot des lacunes.

\subsection{Conjugaison affine des translations et des dilatations}

Les translations et les dilatations satisfont

\[
M_aT_bM_a^{-1}=T_{ab}
\]

lorsque \(a\) est inversible. Cette identité exprime que le produit réorganise les unités de somme. Lorsque \(a\) cesse d’être inversible, la conjugaison se rompt parce que des branches disparaissent. Le radical \(3,6,9\) localise exactement cet échec.


\section{Algèbre des chemins : composition multiplicative, action additive et superposition}

Dans une algèbre des chemins, la somme agrège les alternatives et le produit concatène les trajectoires. Pour des poids locaux \(w(e)\),

\[
K(x,y)
=\sum_{\gamma:x\to y}
\prod_{e\in\gamma}w(e).
\]

Si l’on fixe une branche logarithmique cohérente,

\[
\prod_{e\in\gamma}w(e)
=\exp\!\left(\sum_{e\in\gamma}\log w(e)\right),
\]

et, avec l’action discrète appropriée,

\[
K(x,y)=\sum_{\gamma:x\to y}e^{-S(\gamma)/\hbar}.
\]

L’équation réunit les opérations :

\begin{itemize}
\item le produit compose les pas successifs ;
\item le logarithme accumule l’action locale ;
\item l’exponentielle reconstruit le poids de la route ;
\item la somme agrège toutes les alternatives.
\end{itemize}

L’APP est le support minimal sur lequel somme et produit agissent sur les mêmes histoires. Le TPK ajoute les données qui empêchent la somme sur les chemins de devenir une somme sur des valeurs sans généalogie.

Dans ce cadre, un chiffre est la sortie d’un lecteur ; une valeur est la projection archimédienne d’une section ; un nombre HMT conserve la famille compatible de chemins qui soutient cette valeur ; une particule conserve cette généalogie sous une réalisation spectrale persistante.


\section{Racine interne de \(-1\), structure exponentielle et géométries elliptique, parabolique et hyperbolique}

\subsection{Racine interne de \(-1\) dans la branche finie}

La matrice d’incidence de Paley fournit un opérateur \(A_W\) avec

\[
A_W^2=-I.
\]

Cela fait de \(\mathbb F_3^6\) un espace tridimensionnel sur

\[
\mathbb F_9=\mathbb F_3[\iota]/(\iota^2+1).
\]

La racine de \(-1\) n’est pas adjointe de l’extérieur comme une notation complexe. Elle émerge comme endomorphisme d’un support ternaire déjà construit.

\subsection{Réalisation continue du type quadratique induit}

Dans la branche réelle elliptique, le générateur \(J_{\rm e}\) satisfait

\[
J_{\rm e}^2=-I,
\]

et

\[
e^{tJ_{\rm e}}=\cos t\,I+\sin t\,J_{\rm e}.
\]

En particulier,

\[
e^{\frac\pi2J_{\rm e}}=J_{\rm e},\qquad
e^{\pi J_{\rm e}}=-I,\qquad
e^{2\pi J_{\rm e}}=I.
\]

Ainsi, \(J_{\rm e}\) est le quart de tour interne et l’équation d’Euler est sa clôture au demi-tour.

\subsection{Support opératoriel commun de \(3\) lectures géométriques}

Soit

\[
\mathbb D=\{z\in\mathbb C:|z|<1\}.
\]

Ce même ensemble contient le diamètre réel \((-1,1)\), l’intérieur complexe, la frontière trigonométrique \(S^1\) et la métrique hyperbolique de Poincaré

\[
ds_{\rm P}^2=\frac{4\ell^2|dz|^2}{(1-|z|^2)^2}.
\]

L'application

\[
R_{\pi/2}(z)=iz
\]

préserve \(|z|\) et \(|dz|\). Par conséquent, c’est simultanément :

\begin{itemize}
\item une multiplication complexe par \(\sqrt{-1}\) ;
\item une rotation trigonométrique de \(90^\circ\) ;
\item une isométrie euclidienne du disque ;
\item une isométrie hyperbolique de Poincaré ;
\item un automorphisme de sa frontière circulaire.
\end{itemize}

La coïncidence est littérale : la même application agit sur le même porteur ; c’est le lecteur métrique qui change.

\subsection{Géométries elliptique, parabolique et hyperbolique de l’équation d’Euler étendue}

Les générateurs

\[
J_{+1}^2=-I,\qquad
J_0^2=0,\qquad
J_{-1}^2=I
\]

produisent

\[
e^{tJ_{+1}}=\cos t\,I+\sin t\,J_{+1},
\]

\[
e^{tJ_0}=I+tJ_0,
\]

\[
e^{tJ_{-1}}=\cosh t\,I+\sinh t\,J_{-1}.
\]

La même série exponentielle se sépare en trois régimes selon le carré du générateur : elliptique, parabolique et hyperbolique. La spécialisation

\[
e^{2\log\varphi\,J_{-1}}=F_{\rm av}^{\,2}
\]

intègre l’auto-échelle de Fibonacci. Euler et Fibonacci apparaissent comme des spécialisations d’une famille exponentielle quadratique commune.

\subsection{Porteur commun, opérateur commun et famille exponentielle quadratique}

La fusion géométrique se formule à trois niveaux :

\begin{enumerate}
\item \textbf{porteur commun :} le disque unité contient le diamètre réel, l’intérieur complexe et le bord circulaire ;
\item \textbf{opérateur commun :} \(z\mapsto iz\) est un quart de tour, une multiplication complexe et une isométrie hyperbolique ;
\item \textbf{famille commune :} l’équation exponentielle tritique réunit les branches elliptique, parabolique et hyperbolique.
\end{enumerate}

Le résultat n’identifie pas les métriques euclidienne et hyperbolique. Il affirme que toutes deux sont des réalisations distinctes sur un même support et partagent l’automorphisme issu de la racine interne de \(-I\).


\section{Restriction de la rotation intérieure, holonomie de bord et torsion minimale}

La construction contient deux objets positifs et différenciés :

\begin{enumerate}
\item le générateur interne \(J_{\rm e}\), avec \(J_{\rm e}^2=-I\) ;
\item l’invariant global
\end{enumerate}
   \[
   \Delta_4=
   \frac{\pi-e\log\pi}{270}
   \left(1+\frac{\pi}{729}\right),
   \]
   appelé torsion minimale et utilisé avant la sélection d’espèce.

La thèse de l’auteur ajoute que la dynamique rotationnelle interne, lorsqu’elle est publiée au bord, est absorbée comme torsion minimale. La relation ne doit pas s’écrire comme une simple égalité \(J_{\rm e}=\Delta_4\) : les objets ont des types différents. Sa forme naturelle exige

\[
\text{rotation intérieure}
\xrightarrow{\operatorname{Res}_{\partial}}
\text{holonomie de bord}
\xrightarrow{\operatorname{Def}}
\text{défaut scalaire}.
\]

Soit \(q_\partial\) le lecteur de frontière, \(\nabla_\partial\) la connexion induite et \(e_\partial\) le repère discret. La torsion de bord a pour type

\[
T_\partial=d_{\nabla_\partial}e_\partial.
\]

L’affirmation de l’auteur acquiert une forme probatoire au moyen d’une fonctionnelle

\[
\Lambda_\partial:
\operatorname{Tors}(\partial X)
\longrightarrow\mathbb R
\]

tel que

\[
\Lambda_\partial(T_\partial)=\Delta_4
\]

et d’un théorème qui fasse procéder \(T_\partial\) de la restriction de \(J_{\rm e}\) par les opérateurs de bord du TPK.

La racine interne de \(-I\), le quart de tour, l’équation d’Euler étendue et la formule de \(\Delta_4\) sont démontrés dans leurs domaines. Le passage de la rotation intérieure à la torsion minimale de bord exige le morphisme \(\operatorname{Res}_{\partial}\) et la fonctionnelle \(\Lambda_{\partial}\) déclarés ; il n’est pas identifié par une égalité scalaire.


\section{Triplicité générationnelle, conjugaison et mélange sur l’état enrichi}

La phrase d’Isidor Isaac Rabi peut servir de référence historique, mais elle ne doit pas régir le titre ni constituer une section autonome. La question scientifique est :

\begin{quote}
Quel mécanisme interne distingue trois ordres stables de réalisation fermionique et conserve cette distinction sous couleur, conjugaison et mélange ?
\end{quote}

La réponse HMT n’est pas « parce que TRIT a trois valeurs ». La structure sépare

\[
\operatorname{flav}(X)
=\bigl(\sigma_{ud}(X),g(X),\chi_{\rm fina}(X)\bigr),
\]

où \(\sigma_{ud}\) distingue les branches, \(g\) distingue la génération physique et \(\chi_{\rm fina}\) conserve l’orientation et la mémoire au sein d’une branche. Dans les quarks, l’orbite de couleur est ajoutée après la saveur et avant le caractère.

La triplicité est obtenue comme conséquence de la structure enrichie ; ses effets sur la couleur, la conjugaison et le mélange sont développés ci-dessous :

\begin{enumerate}
\item couleur, saveur, génération et masse sont des lecteurs distincts du même état enrichi ;
\item la conjugaison particule–antiparticule agit sur l’orientation sans effacer la classe générationnelle ;
\item CKM et PMNS transportent des bases d’états déjà typés ; ils ne créent pas la triplicité ;
\item la classification neutrinique doit découler de la représentation, de l’occupation et du mélange, non d’un changement de nom des profondeurs \(G0,\ldots,G4\).
\end{enumerate}

Le fait que les neutrinos soient des leptons de spin demi-entier et d’obéissance fermionique est une classification physique connue. L’apport structurel HMT consiste à obtenir intérieurement pourquoi le secteur neutre occupe cette représentation, pourquoi il admet trois réalisations physiques et comment il se couple à l’opérateur de mélange sans lire ses angles externes.


\section{Mémoire opératorielle du TPK : conservation du résidu, du quotient, de la retenue, du feuillet, de l’orientation, de la route et de la frontière sous projection}

\subsection{Perte de réversibilité et relèvement généalogique}

Dans l’APP, différentes histoires peuvent produire le même résidu. Le produit replie davantage d’histoires que la somme ; la puissance en replie davantage que le produit ; la racine tente de reconstruire des branches déjà identifiées. Le TPK conserve la donnée qui permet d’inverser cette perte :

\[
\text{résidu},\quad
\text{quotient},\quad
\text{retenue},\quad
\text{feuillet},\quad
\text{orientation},\quad
\text{route},\quad
\text{frontière},\quad
\text{mémoire}.
\]

L’opération algébrique peut être caractérisée non seulement par sa valeur, mais par le type de généalogie qu’elle conserve ou détruit.

\subsection{Projection des histoires compatibles sur le continu archimédien}

La droite réelle apparaît lorsqu’une famille de cylindres compatibles se contracte sous l’évaluation archimédienne. La valeur réelle est l’ombre de la section ; non sa généalogie complète. La projection exceptionnelle conserve une autre partie comme incidence et symétrie.

La somme sur les chemins exprime le même principe :

\[
\text{alternatives}\xrightarrow{\sum}\text{amplitude globale},
\qquad
\text{étapes successives}\xrightarrow{\prod}\text{poids de route}.
\]

Le logarithme transforme la composition multiplicative en action additive ; l’exponentielle la ramène à une amplitude. La structure discrète du continu peut être formulée comme conservation généalogique de toutes les routes avant leur projection sur \(\mathbb R\).

\subsection{Distinction entre section numérique et réalisation persistante de particule}

Un nombre et une particule partagent des invariants ontologiques minimaux : mémoire, orientation, persistance, clôture, signature et position dans une famille de prolongements.

Le nombre publie une coordonnée archimédienne. La particule ajoute une réalisation spectrale, une occupation, une propagation et une stabilité. Le passage HMT \(\to\) MD change le lecteur appliqué à une même généalogie enrichie ; il ne change pas arbitrairement d’objet.


\section{Clôture conjointe du réseau radical, de la loi puissance--lacune, de la géométrie d’Euler et de la hiérarchie opératorielle}

\subsection{Théorème du réseau radical nonadique}

Le réseau \(3,6,9\) de l’APP est la représentation visuelle du radical nilpotent \((3)\subset\mathbb Z/9\mathbb Z\). Ses phases \(369\), \(693\) et \(936\) réalisent l’isomorphisme \(R/(3)\cong(3)\) induit par la multiplication par trois. Les bandes localisent exactement le défaut d’inversibilité du produit.

\subsection{Théorème puissance--lacune}

Pour tout \(n\ge1\), le nombre de rangs décimaux perdus par \(9^n\) par rapport à \(10^n\) est \(\lceil n\log_{10}(10/9)\rceil\). Ses accroissements forment le mot mécanique des lacunes. En \(n=9^9\), la coordonnée retenue est \(369\,693\,099\) et la longueur résultante est \(369\,693\,100\).

\subsection{Théorème géométrique de la famille exponentielle d’Euler}

La racine interne de \(-I\) fournit un quart de tour. Dans le disque unité, la même application \(z\mapsto iz\) est une multiplication complexe, une rotation trigonométrique et une isométrie de Poincaré. L’équation exponentielle tritique prolonge ce régime aux régimes parabolique et hyperbolique, en incorporant dans le dernier l’auto-échelle de Fibonacci.

\subsection{Principe de conservation opératorielle entre niveaux}

La hiérarchie somme–produit–puissance est une hiérarchie de fibres et de perte de réversibilité. Le TPK est le relèvement qui conserve les coordonnées nécessaires pour distinguer les histoires identifiées par les lecteurs ordinaires.


\section{Hypothèses de typage et contrôles négatifs}

\begin{enumerate}
\item Ne pas identifier \(369\,693\,100\) à un mot TPK sans déclarer le lecteur entre les domaines.
\item Ne pas confondre l’idéal \(\mathfrak m^2=0\) avec \(\mathfrak m\times\mathfrak m\).
\item Ne pas confondre modulo neuf et racine numérique : \(9\mapsto0\) dans le premier et \(9\mapsto9\) dans la seconde.
\item Ne pas utiliser la triplicité de TRIT comme démonstration suffisante de trois générations physiques.
\item Ne pas transformer \(J^2=-I\) en \(\Delta_4\) sans l’application de bord et la fonctionnelle de torsion.
\item Ne pas présenter \(369\) comme un préfixe de \(9^{9^9}\) : il appartient à son ordre et à sa longueur décimale.
\item Ne pas réduire les dix mille chiffres à une preuve de grandeur : ils certifient des troncatures étendues d’un processus paramétrable dont la thèse est le prolongement indéfini.
\end{enumerate}

\section{Conséquences conjointes pour le réseau radical, la mémoire opératorielle et la torsion de bord}

La découverte principale n’est pas qu’une tour contienne visuellement les chiffres \(3,6,9\). La puissance nonadique et le calendrier des lacunes sont deux présentations du même défaut logarithmique, tandis que l’APP localise \(3,6,9\) comme le radical où la multiplication perd son inversibilité. La tour relie les deux extrémités :

\[
\text{radical nonadique}
\longrightarrow
\text{puissance}
\longrightarrow
\text{logarithme}
\longrightarrow
\text{vacance}
\longrightarrow
\text{ordre décimal}
\longrightarrow
\text{clôture }9\to1.
\]

L’équation d’Euler étendue ajoute la deuxième direction : la racine interne de \(-I\) transforme la structure discrète en dynamique de rotation et réunit, sur un support commun, la droite réelle, le cercle trigonométrique et le disque hyperbolique. Le passage ultérieur à la torsion minimale doit conserver l’application de bord déjà développée dans le corpus, non la remplacer par une identification nominale.
